\subsubsection{Definições}

Algumas definições e termos importantes:

\paragraph{Calcular teto}

Para calcular o teto da divisão entre a e b:
\begin{center}
a + b - 1 / b
\end{center}

\paragraph{Funções}

\begin{itemize}\setlength\itemsep{0pt}
    \item \textbf{Comutativa:} Uma função $f(x, y)$ é comutativa se $f(x, y) = f(y, x)$.
    \item \textbf{Associativa:} Uma função $f(x, y)$ é associativa se $f(x, f(y, z)) = f(f(x, y), z)$.
    \item \textbf{Idempotente:} Uma função $f(x, y)$ é idempotente se $f(x, x) = x$.
\end{itemize}

\subsubsection{Números primos}

Números primos são muito úteis para funções de hashing (dentre outras coisas). Aqui vão vários números primos interessantes:

\paragraph{Primos com truncamento à esquerda}

Números primos tais que qualquer sufixo deles é um número primo:
\begin{center}
33,333,31\\
357,686,312,646,216,567,629,137
\end{center}

\paragraph{Primos gêmeos (Twin Primes)}

Pares de primos da forma $(p, p+2)$ (aqui tem só alguns pares aleatórios, existem muitos outros).

\paragraph{Números primos de Mersenne}

São os números primos da forma $2^m - 1$, onde $m$ é um número inteiro positivo.

\subsubsection{Operadores lineares}

\paragraph{Rotação no sentido anti-horário por $\theta$}

\[
\begin{bmatrix}
\cos \theta & -\sin \theta \\
\sin \theta & \cos \theta
\end{bmatrix}
\]

\paragraph{Reflexão em relação à reta $y = mx$}

\[
\frac{1}{m^2 + 1}
\begin{bmatrix}
1 - m^2 & 2m \\
2m & m^2 - 1
\end{bmatrix}
\]

\paragraph{Inversa de uma matriz $2 \times 2$ $A$}

\[
\begin{bmatrix}
a & b \\
c & d
\end{bmatrix}^{-1}
= \frac{1}{\det(A)}
\begin{bmatrix}
d & -b \\
-c & a
\end{bmatrix}
\]

\paragraph{Cisalhamento horizontal por $K$}

\[
\begin{bmatrix}
1 & K \\
0 & 1
\end{bmatrix}
\]

\paragraph{Cisalhamento vertical por $K$}

\[
\begin{bmatrix}
1 & 0 \\
K & 1
\end{bmatrix}
\]

\paragraph{Mudança de base}

$\vec{a}_\beta$ são as coordenadas do vetor $\vec{a}$ na base $\beta$.\\
$\vec{a}$ são as coordenadas do vetor $\vec{a}$ na base canônica.\\
$\vec{b}_1$ e $\vec{b}_2$ são os vetores de base para $\beta$.\\
$C$ é uma matriz que muda da base $\beta$ para a base canônica.

\[
C\vec{a}_\beta = \vec{a}
\]
\[
C^{-1}\vec{a} = \vec{a}_\beta
\]
\[
C = \begin{bmatrix}
{b_1}_x & {b_2}_x \\
{b_1}_y & {b_2}_y
\end{bmatrix}
\]

\paragraph{Propriedades das operações de matriz}

\begin{align*}
(AB)^{-1} &= B^{-1}A^{-1} \\
(AB)^T &= B^T A^T \\
(A^{-1})^T &= (A^T)^{-1} \\
(A + B)^T &= A^T + B^T \\
\det(A) &= \det(A^T) \\
\det(AB) &= \det(A)\det(B)
\end{align*}

Seja $A$ uma matriz $N \times N$:
\[
\det(kA) = k^N \det(A)
\]

\subsubsection{Sequências numéricas}

\paragraph{Sequência de Fibonacci}

Primeiros termos: 1, 1, 2, 3, 5, 8, 13, 21, 34, \dots

Definição:
\begin{align*}
F_0 &= F_1 = 1 \\
F_n &= F_{n-1} + F_{n-2}
\end{align*}

Matriz de recorrência:
\[
\begin{bmatrix}
0 & 1 \\
1 & 1
\end{bmatrix}
\begin{bmatrix}
F_{n-2} \\
F_{n-1}
\end{bmatrix}
=
\begin{bmatrix}
F_{n-1} \\
F_n
\end{bmatrix}
\]

\paragraph{Sequência de Catalan}

Primeiros termos: 1, 1, 2, 5, 14, 42, 132, 429, 1430, \dots

Definição:
\begin{align*}
C_0 &= C_1 = 1 \\
C_n &= \sum_{i=0}^{n-1} C_i \cdot C_{n-1-i}
\end{align*}

Definição analítica:
\[
C_n = \frac{1}{n+1} \binom{2n}{n}
\]

Propriedades úteis:
\begin{itemize}\setlength\itemsep{0pt}
    \item $C_n$ é o número de árvores binárias com $n + 1$ folhas.
    \item $C_n$ é o número de sequências de parênteses bem formadas com $n$ pares de parênteses.
\end{itemize}

\subsubsection{Análise combinatória}

\paragraph{Fatorial}

O fatorial de um número $n$ é o produto de todos os inteiros positivos menores ou iguais a $n$.\\
O fatorial conta o número de permutações de $n$ elementos.
\[
n! = n \cdot (n - 1)!
\]
Em particular, $0! = 1$.

\paragraph{Combinação}

Conta o número de maneiras de escolher $k$ elementos de um conjunto de $n$ elementos.
\[
\binom{n}{k} = \frac{n!}{k! \cdot (n - k)!}
\]

\paragraph{Arranjo}

Conta o número de maneiras de escolher $k$ elementos de um conjunto de $n$ elementos, onde a ordem importa.
\[
P(n, k) = \frac{n!}{(n - k)!}
\]

\paragraph{Estrelas e barras}

Conta o número de maneiras de distribuir $n$ elementos idênticos em $k$ recipientes distintos.
\[
\binom{n + k - 1}{k - 1}
\]

\paragraph{Princípio da inclusão-exclusão}

O princípio da inclusão-exclusão é uma técnica para contar o número de elementos em uma união de conjuntos.
\[
\left| \bigcup_{i=1}^n A_i \right| = \sum_{k=1}^n (-1)^{k+1} \sum_{1 \le i_1 < \dots < i_k \le n} |A_{i_1} \cap \dots \cap A_{i_k}|
\]

\paragraph{Princípio da casa dos pombos}

Se $n$ pombos são colocados em $m$ casas, então pelo menos uma casa terá $\left\lceil \frac{n}{m} \right\rceil$ pombos ou mais.

\subsubsection{Teoria dos números}

\paragraph{Pequeno teorema de Fermat}

Se $p$ é um número primo e $a$ é um inteiro não divisível por $p$, então $a^{p-1} \equiv 1 \pmod p$.

\paragraph{Teorema de Euler}

Se $m$ e $a$ são inteiros positivos coprimos, então $a^{\phi(m)} \equiv 1 \pmod m$, onde $\phi(m)$ é a função totiente de Euler.

\paragraph{Aritmética modular}

Quando estamos trabalhando com aritmética módulo um número $p$, todos os valores existentes estão entre $[0, p-1]$.\\
Algumas propriedades e equivalências úteis para usar aritmética modular em código são:
\begin{itemize}\setlength\itemsep{0pt}
    \item $(a + b) \bmod p \equiv ((a \bmod p) + (b \bmod p)) \bmod p$
    \item $(a - b) \bmod p \equiv ((a \bmod p) - (b \bmod p)) \bmod p$
    \begin{itemize}\setlength\itemsep{0pt}
        \item Note que o resultado pode ser negativo, nesse caso é necessário adicionar $p$ ao resultado. De forma geral, geralmente fazemos $(a - b + p) \bmod p$ (assumindo que $a$ e $b$ já estão no intervalo $[0, p-1]$).
    \end{itemize}
    \item $(a \cdot b) \bmod p \equiv ((a \bmod p) \cdot (b \bmod p)) \bmod p$
    \item $a^b \bmod p \equiv ((a \bmod p)^b) \bmod p$
    \item $a^b \bmod p \equiv ((a \bmod p)^{(b \bmod \phi(p))}) \bmod p$ se $a$ e $p$ são coprimos
\end{itemize}